% Editar este fichero y rellenar con el resumen del TFM en inglés
Accurate blood glucose prediction is a key tool for type 1 diabetes management, enabling anticipation of hypoglycemic and hyperglycemic episodes and improving patient quality of life. Deep learning models based on Long Short-Term Memory (LSTM) networks have shown strong performance in this domain due to their ability to capture complex temporal dependencies in continuous glucose monitoring (CGM) signals. However, the inherent demographic heterogeneity of patient cohorts — particularly differences in age and sex — can introduce systematic biases during training, undermining the fairness and generalisability of the resulting models.

This Master's thesis addresses this problem through the design, implementation, and evaluation of a data balancing pipeline applied exclusively to the training set, within a Group K-Fold cross-validation scheme that enforces strict patient separation across all folds. The central objective is to quantify whether demographic balancing techniques — based on age and sex — reduce patient heterogeneity bias and improve the clinical performance of LSTM models, measured using glycaemic time-in-range metrics (Time in Range, TIR; Time Above Range, TAR; Time Below Range, TBR).

Three CGM datasets are used: \textbf{T1DiabetesGranada} (643 patients, over 30 million readings collected between 2018 and 2022 within the Granada healthcare system), \textbf{REPLACE-BG} (a US multicentre clinical trial), and \textbf{DiaTrend} (an international open-access cohort). This diversity allows the robustness of findings to be assessed across datasets with distinct structural characteristics — size, density, demographic distribution, and missing data patterns.

The implemented technique taxonomy comprises nine experimental conditions organised into three categories: random techniques (random undersampling, random oversampling), guided techniques (SMOTE, Tomek Links), and hybrid combinations (patient-aware undersampling, jittering, undersampling+oversampling, undersampling+SMOTE, SMOTE+Tomek), each applied in two demographic variants (age-based and sex-based), plus an unbalanced control condition. The pipeline was formally validated through a constraint-verification notebook checking, among other properties, the bit-identical invariance of validation and test sets across conditions, the absence of patient overlap, and the compliance of SMOTE-generated synthetic points with sensor measurement ranges.

The statistical analysis is structured using Friedman tests to detect global differences across conditions, Nemenyi post-hoc tests for pairwise comparisons, Cohen's $d$ effect sizes, and Borda count ranking aggregation. Results reveal a dataset-dependent picture: \textbf{DiaTrend} is the only dataset in which balancing yields statistically significant effects across all analysed glycaemic ranges, while \textbf{T1DiabetesGranada} and \textbf{REPLACE-BG} exhibit low fold-level ranking consistency — measured via Spearman correlation — attributable to intrinsic cohort heterogeneity and the Missing Not At Random (MNAR) missingness pattern identified during the exploratory phase.

Regarding demographic dimensions, \textbf{age-based} balancing consistently outperforms \textbf{sex-based} balancing in terms of clinical impact. \textbf{Age·Oversampling} achieves the greatest TIR improvement, albeit at the cost of systematic TAR-1 degradation, revealing a clinical trade-off that must be carefully considered in practice. Global Borda aggregation positions \textbf{Age·SMOTE} as the best-performing technique overall, offering a favourable balance between improvement in the euglycaemic range and control of adverse effects in extreme ranges.

This work therefore contributes to the intersection of algorithmic fairness and clinical time-series modelling, providing empirical evidence on when and to what extent demographic balancing is beneficial, and laying the methodological groundwork for its incorporation into glucose prediction systems designed for heterogeneous patient populations.